\documentclass[10pt,a4paper]{amsart}
\usepackage[margin=19mm]{geometry}
\usepackage{amsmath,amssymb,mathtools}
\usepackage[scr=rsfs,cal=euler]{mathalfa}
\usepackage{xfrac}
\usepackage[unicode,hidelinks,bookmarks=false]{hyperref}
\newcommand{\ie}{i.e.\ }
\newcommand{\cf}{cf.\ }
\newcommand{\resp}{respectively\ }
\newcommand{\Subsection}{\subsection}
\providecommand{\nobreakdash}{\nobreak\hbox{-}}
\setlength{\emergencystretch}{2em}
\multlinegap0pt
\usepackage{bm}
\usepackage{bbm}
\usepackage{dsfont}
\usepackage{xspace}
\usepackage{cancel}
\usepackage{mathtools}
\usepackage{amsthm}
\makeatletter
\@ifundefined{theorem}{\newtheorem{theorem}{Theorem}}{}
\@ifundefined{lemma}{\newtheorem{lemma}[theorem]{Lemma}}{}
\@ifundefined{remark}{\newtheorem{remark}[theorem]{Remark}}{}
\makeatother
\makeatletter
\newcommand{\R}{\mathbb R}
\expandafter\ifx\csname remark\endcsname\relax
\newenvironment{remark}{\begin{preremark}\rm}{\medskip \end{preremark}}
\fi
\makeatother
\title[Online Koml\'os converges to mean curvature flow]{Online Koml\'os converges to mean curvature flow}
\author{Nestor Guillen, Vladimir A. Kobzar}
\date{Source: arXiv:2607.08943v1 (CC BY 4.0)}
\begin{document}
\maketitle

\noindent\emph{Source and licence.} Online Koml\'os converges to mean curvature flow. Version arXiv:2607.08943v1 (CC BY 4.0), distributed under the Creative Commons Attribution 4.0 International licence, as recorded in the arXiv licence metadata for this exact version and shown on the abstract page https://arxiv.org/abs/2607.08943v1. Full rights notice: licenses/arxiv-2607.08943-CC-BY-4.0.txt.\par\smallskip

\noindent\textit{Editorial scope.} Counted unit: the Remark whose source label is \texttt{r:half relaxed limits H n delta} in the article (source0.tex lines 842--850, source label \texttt{r:half relaxed limits H n delta}), delivered together with the article's own complete original proof of it (source0.tex lines 852--950, one \texttt{proof} environment of 99 lines). No mathematics is written, repaired or re-derived: statement and proof are the article's own text line for line. Required context retained uncounted: e:definition of the auxiliary operator H delta n (lines 781--783); l:liminf and limsup for H delta n (lines 824--840). Every definition, display and label of those lines is retained unchanged. Omitted from this package and not redistributed: 1--780, 784--823, 841--841, 851--851, 951--1800. Nothing inside a retained excerpt is omitted or altered: every retained body is delivered exactly as the article's lines are written, so no omission receipt of any kind accompanies this package. Citations: the retained excerpts carry no citation command at all, so no bibliography entry of the article is transcribed and no reference is redistributed. Reference adaptation: none was needed and none was made; every reference inside the delivered unit resolves inside it, so no unresolved reference or citation is produced. Printed slips kept literal: no printed word, symbol, hypothesis or number of the counted statement or of its proof was altered. Layout and class: the article's own class is \texttt{amsart} with packages including amsmath, amsfonts, amssymb, amsthm, esint, mathtools, geometry, graphicx, xcolor, cite, hyperref, comment, color, dsfont; the wrapper is the lane's pinned \texttt{amsart} editorial wrapper at 10pt on A4, which restates the theorem-like environments the retained text needs and adds the article's own macro definitions that the retained text needs (\texttt{\textbackslash R}), with extra wrapper packages (bm, bbm, dsfont, xspace, cancel, mathtools). The delivered numbering is the unit's own. Assets: no retained span contains a figure environment, a graphics-inclusion command, a PDF-page inclusion, a table or a TikZ picture; no image file is copied into this package, the package declares no asset dependency and no image file is redistributed with it. Licence: version arXiv:2607.08943v1 of this article is distributed under the Creative Commons Attribution 4.0 International licence, as recorded in the arXiv licence metadata for this exact version and shown on the abstract page https://arxiv.org/abs/2607.08943v1. A full rights notice is delivered at licenses/arxiv-2607.08943-CC-BY-4.0.txt. Characters: every retained excerpt is pure ASCII, so the delivered engine, XeLaTeX with its default Latin Modern text font, reports no missing character.
\begin{align}\label{e:definition of the auxiliary operator H delta n}
  H_{\delta,n}(M,p) := \max_A\min_{\overline\varepsilon}\left \{ 2\delta^{-1}(p,A\overline\varepsilon) + (M A\overline\varepsilon,A\overline\varepsilon) \right \}.
\end{align}

\begin{lemma}\label{l:liminf and limsup for H delta n}
For any sequences
\begin{align*}
  M_k \to M,\; p_k \to p,\; \delta_k \to 0,    \end{align*}
  of symmetric matrices $M_k \in \R^{m\times m}$, vectors $p_k \in \R^m$ and scalars $\delta_k>0$,
we have the limit
\begin{align*}
  \lim \limits_{k\to \infty} H_{\delta_k,n}(M_k,p_k) = H_n(M,p),\;\text{ when } p \neq 0,
\end{align*}
as well as the bounds
\begin{align*}
  \liminf \limits_{k\to \infty}H_{\delta_k,n}(M_k,p_k) & \geq (H_n)_*(M,0)~~~ \text { and} ~~~
  \limsup \limits_{k\to \infty}H_{\delta_k,n}(M_k,p_k)  \leq (H_n)^*(M,0) \text{ when } p=0.
\end{align*}
Moreover, both bounds are sharp, i.e., for both bounds one can find sequences where the bounds hold with equality.

\end{lemma}

\begin{remark}\label{r:half relaxed limits H n delta}
  Another way of stating Lemma \ref{l:liminf and limsup for H delta n} is in terms of the half-relaxed limits of $H_{\delta,n}$ as $\delta \to 0$. In particular, the lemma says that when $p =0$,
  \begin{align*}
    (H_n)^*(M,p) & = \limsup_{M'\to M, p'\to p,\delta \to 0} H_{\delta,n}(M',p') = \max_A \min_{\overline \varepsilon}(MA\overline\varepsilon,A\overline\varepsilon),\\
    (H_n)_*(M,p) & = \liminf_{M'\to M, p'\to p,\delta \to 0} H_{\delta,n}(M',p') = \min_{\sigma \in \mathbb{S}^{m-1}} H_n(M,\sigma), 
  \end{align*}
  and furthermore, $(H_n)^*(M,p) = (H_n)_*(M,p) = H_n(M,p)$ whenever $p \neq 0$. 

\end{remark}

\begin{proof} We start by fixing an arbitrary $\delta>0$ and recording two observations about $H_{\delta,n}$. First, it is clear that restricting the maximization in $A$ to just those $A$'s with $A\in \mathcal{A}(p)$ cannot increase the maximum, and so we have for any symmetric $M$ and vector $p$,
\begin{align}\label{e:H delta n basic lower pointwise bound}
  H_{\delta,n}(M,p) \geq \max_{A\in\mathcal{A}(p)}\min_{\overline\varepsilon}(MA\overline\varepsilon,A\overline\varepsilon) = H_n(M,p).    
\end{align}
Here, we have made use of the fact that $A^\top p = 0$ whenever $A \in \mathcal{A}(p)$. 

The second observation concerns $H_{\delta,n}$ when $p\neq 0$, we are going to show that given any symmetric $M$ and a vector $p\neq 0$, if $A$ achieves the maximum in \eqref{e:definition of the auxiliary operator H delta n} we have 
\begin{align}\label{e:H delta n optimal A bound for non zero p}
  \|(I-\Pi(p))A\| & = 
\|\hat p \otimes \hat p A\| \leq \tfrac{1}{2}\delta |p|^{-1} C(m,n)\|M\|,  
\end{align}
where $\Pi(p) := I-\hat p\otimes \hat p$, $\hat p := p/|p|$ for $p\neq 0$. To prove \eqref{e:H delta n optimal A bound for non zero p} we first note that the minimization in \eqref{e:definition of the auxiliary operator H delta n} involves an objective function that is the sum of an odd and an even function of $\overline \varepsilon$. Keeping this in mind, we compare the values of the objective function at $\overline\varepsilon$ and $-\overline \varepsilon$, and conclude that
\begin{align}
  \min_{\overline\varepsilon}\{2\delta^{-1}(p,A\overline \varepsilon)+ (MA\overline \varepsilon,A\overline \varepsilon)\} & =   \min_{\overline\varepsilon}\{ -2\delta^{-1}|(p,A\overline \varepsilon)|+ (MA\overline \varepsilon,A\overline \varepsilon)\}. \label{e:odd_even_relationship}
\end{align}
Then, writing $A = \Pi(p)A+(I-\Pi(p))A$, where , 
\begin{align*}
  (p,A\overline \varepsilon) = (p,(I-\Pi(p))A\overline \varepsilon) =  |p|(A^\top \hat p,\overline\varepsilon). 
\end{align*}
Then, given that $A$ achieves the maximum  in \eqref{e:definition of the auxiliary operator H delta n} for a given $M$ and $p$,  
\begin{align*}
 0 \leq \min_{\overline \varepsilon}\{ -2\delta^{-1}|(p,A\overline \varepsilon)|+ (MA\overline \varepsilon,A\overline \varepsilon)\},
\end{align*}
and thus for every $\overline \varepsilon$ we have
\begin{align*}
  2\delta^{-1}|(A^\top p,\overline \varepsilon)| \leq (MA\overline \varepsilon, A\overline \varepsilon) \leq C(m,n)\|M\|.
\end{align*}
Then, choosing $\overline \varepsilon$ for which
\begin{align*}
  (A^\top p,\overline \varepsilon)= |(Ae_1)^\top\cdot p| + \ldots + |(Ae_n)^\top\cdot p| = \|A^\top p\|_1,   
\end{align*}
we conclude that 
\begin{align*}
  \|A^\top p\|_1 \leq \tfrac{1}{2}\delta C(m,n)\|M\|.      
\end{align*}
Then, noting that $\|\hat p \otimes \hat p A\| =| A^\top \hat p| \leq |p|^{-1}\| A^\top  p\|_1$ we obtain \eqref{e:H delta n optimal A bound for non zero p}.

With these basic estimates in hand, consider a generic sequence $M_k,p_k,\delta_k$ as in the statement of the lemma. First, take the case where $p_k \to p \neq 0$. If we let $A_k$ denote a matrix achieving the maximum for $(M_k,p_k)$, estimate \eqref{e:H delta n optimal A bound for non zero p} says (without loss of generality, $p_k \neq 0$ for every $k$)
\begin{align*}
  \|(I-\Pi(p_k))A_k\|\leq \tfrac{1}{2}\delta_k |p_k|^{-1} C(m,n)\|M_k\|\;\;\forall\;k.
\end{align*}
On account of $p \neq 0$ we know that for all sufficiently large $k$ we have $\tfrac{1}{2}|p_k|^{-1}\leq |p|^{-1}$. Likewise, for large enough $k$ we have $\|M_k\|\leq 1+ 2\|M\|$. Therefore for all large enough $k$ we have the inequality
\begin{align*}
  \|(I-\Pi(p_k))A_k\| \leq \delta_k |p|^{-1}C(m,n)(1+ 2\|M\|).    
\end{align*}
We use this to estimate $(M_kA_k\overline \varepsilon,\overline \varepsilon)$ for large $k$, via the decomposition $A_k = \Pi(p_k)A_k+(I-\Pi(p_k))A_k$,  
\begin{align*}
  (M_kA_k\overline \varepsilon, A_k\overline \varepsilon) \leq (M_k\Pi(p_k)A_k \overline\varepsilon,\Pi(p_k)A_k \overline \varepsilon) + C\|M_k\|\|(I-\Pi(p_k))A_k\|. 
\end{align*}
Then for all large enough $k$ and for every $\overline \varepsilon$ we have
\begin{align*}
  (M_kA_k\overline \varepsilon, A_k\overline \varepsilon) \leq \left \{ (M_k\Pi(p_k)A_k \overline\varepsilon,\Pi(p_k)A_k \overline \varepsilon)\right \} + C(m,n)(1+2\|M\|)^2|p|^{-1}\delta_k. 
\end{align*}
 From here it follows that
\begin{align*}
   H_{\delta_k,n}(M_k,p_k) & \leq \min_{\overline\varepsilon}(M_k(\Pi(p_k)A_k)\overline\varepsilon, (\Pi(p_k)A_k)\overline \varepsilon) + C(m,n)(1+2\|M\|)^2|p|^{-1}\delta_k\\
   & \leq H_n(M_k,p_k) + C(m,n)(1+2\|M\|)^2|p|^{-1}\delta_k,
\end{align*}
where we used that $\Pi(p_k)A_k \in \mathcal{A}(p_k)$ to obtain the second inequality. Passing to the limit, we have
\begin{align*}
   \limsup_{k \to \infty} H_{\delta_k,n}(M_k,p_k) & \leq H_n(M,p).
\end{align*}
The case $p\neq 0$ now follows simply by applying \eqref{e:H delta n basic lower pointwise bound} to $H_{\delta_k}(M_k,p_k)$, which leads to 
\begin{align*}
   \liminf \limits_{k \to \infty} H_{\delta_k,n}(M_k,p_k) & \geq \lim \limits_{k\to \infty} H_n(M_k,p_k) = H_n(M,p).
\end{align*}
This shows $(H_n)^*(M,p) = (H_n)_*(M,p) = H_n(M,p)$ for $p\neq 0$.

Now we consider what happens when $p_k \to 0$. For the liminf bound, we  use equation \eqref{e:H delta n basic lower pointwise bound} once again:
\begin{align*}
  H_{\delta,n}(M,p) \geq H_{n}(M,p) \geq \min_{q}H_{n}(M,q).
\end{align*}
The minimum being over all $q\in\mathbb{R}^m$, noting that the minimum cannot be achieved only at $q=0$, and using that $H_n(M,q)$ is $0$-homogeneous for $q\neq 0$, we obtain in fact that  
\begin{align*}
  H_{\delta,n}(M,p) \geq \min_{\sigma \in \mathbb{S}^{m-1}}H_{n}(M,\sigma).
\end{align*}
Since this is true for every $M,p,\delta$, it follows that
\begin{align*}
  \liminf_{k \to \infty} H_{\delta_k,n}(M_k,p_k) \geq \min_{\sigma \in \mathbb{S}^{m-1}}H_n(M,\sigma).    
\end{align*}
Now for the limsup bound, we use equation \eqref {e:odd_even_relationship}, which implies that
\begin{align*}
  \max_A \min_{\overline\varepsilon}\{2\delta^{-1}(p,A\overline \varepsilon)+ (MA\overline \varepsilon,A\overline \varepsilon)\} \leq \max_A \min_{\overline\varepsilon} (MA\overline \varepsilon,A\overline \varepsilon) = H_n(M,0).
\end{align*}
Then, considering the sequence $M_k,p_k,\delta_k$ we have
\begin{align*}
  \limsup_{k \to \infty} H_{\delta_k,n}(M_k,p_k) \leq H_n(M,0),
\end{align*}
which proves the bound for the limsup. To see the limsup upper bound is optimal, simply consider a sequence $M_k,p_k,\delta_k$ where $p_k = 0$ for every $k$. To see why the lower bound for the liminf is optimal, fix an arbitrary $\sigma_0 \in \mathbb{S}^{m-1}$ and consider a sequence with $p_k = (\delta_k)^{1/2}\sigma_0$, in this case, we have
\begin{align*}
  \lim_{k \to \infty} H_{\delta_k,n}(M_k,p_k) = H_n(M,\sigma_0).  
\end{align*}
Indeed, this follows by applying \eqref{e:H delta n optimal A bound for non zero p}, and noting that $\delta_k|p_k|^{-1} = \delta_k^{1/2}$, in which case the argument used for the case $p_k \to p \neq 0$ says that
\begin{align*}
  \|(I-\Pi(p_k))A_k\| \leq \delta_k^{1/2}(1+2\|M\|)    
\end{align*}
and so the distance of $A_k$ to $\mathcal{A}(\sigma_0)$ goes to zero as $k\to \infty$, proving the last limit holds. Since $\sigma_0$ is arbitrary, we can choose it to be the minimizer of $\sigma \mapsto H_n(M,\sigma)$, and this completes the proof.

\end{proof}

\end{document}
